\section{AIME II 2014}
        \begin{enumerate}
        	\item (\textit{AIME II 2014}) Abe can paint the room in 15 hours, Bea can paint 50 percent faster than Abe, and Coe can paint twice as fast as Abe. Abe begins to paint the room and works alone for the first hour and a half. Then Bea joins Abe, and they work together until half the room is painted. Then Coe joins Abe and Bea, and they work together until the entire room is painted. Find the number of minutes after Abe begins for the three of them to finish painting the room.


 
 \textbf{Answer: }334



\textbf{Solution}

From the given information, we can see that Abe can paint $\frac{1}{15}$ of the room in an hour, Bea can paint $\frac{1}{15}\times\frac{3}{2} = \frac{1}{10}$ of the room in an hour, and Coe can paint the room in $\frac{1}{15}\times 2 = \frac{2}{15}$ of the room in an hour.  After $90$ minutes, Abe has painted $\frac{1}{15}\times\frac{3}{2}=\frac{1}{10}$ of the room.  Working together, Abe and Bea can paint $\frac{1}{15}+\frac{1}{10}=\frac{1}{6}$ of the room in an hour, so it takes then $\frac{2}{5}\div \frac{1}{6}= \frac{12}{5}$ hours to finish the first half of the room.  All three working together can paint $\frac{1}{6}+\frac{2}{15}=\frac{3}{10}$ of the room in an hour, and it takes them $\frac{1}{2}\div \frac{3}{10}=\frac{5}{3}$ hours to finish the room.  The total amount of time they take is
 \[90+\frac{12}{5}\times 60+\frac{5}{3}\times 60 = 90+ 144 + 100 = \boxed{334} \text{\ minutes.}\]

 




	\item (\textit{AIME II 2014}) Arnold is studying the prevalence of three health risk factors, denoted by A, B, and C, within a population of men. For each of the three factors, the probability that a randomly selected man in the population has only this risk factor (and none of the others) is 0.1. For any two of the three factors, the probability that a randomly selected man has exactly these two risk factors (but not the third) is 0.14. The probability that a randomly selected man has all three risk factors, given that he has A and B is $\frac{1}{3}$. The probability that a man has none of the three risk factors given that he does not have risk factor A is $\frac{p}{q}$, where $p$ and $q$ are relatively prime positive integers. Find $p+q$. 


 
 \textbf{Answer: }076



\textbf{Solution}

We first assume a population of $100$ to facilitate solving. Then we simply organize the statistics given into a Venn diagram.

 
                    \begin{figure}[H]
                        \centering
                        \includegraphics[width=0.5\textwidth]{Solution/AIME_2014_II_2_Solution_Figure_1.png}
                    \end{figure}
                    

 Let $x$ be the number of men with all three risk factors.  Since "the probability that a randomly selected man has all three risk factors, given that he has A and B is $\frac{1}{3}$," we can tell that $x = \frac{1}{3}(x+14)$, since there are $x$ people with all three factors and 14 with only A and B.  Thus $x=7$.

 Let $y$ be the number of men with no risk factors.  It now follows that
 \[y= 100 - 3 \cdot 10 - 3 \cdot 14 - 7 = 21.\]The number of men with risk factor A is $10+2 \cdot 14+7 = 45$ (10 with only A, 28 with A and one of the others, and 7 with all three).  Thus the number of men without risk factor $A$ is 55, so the desired conditional probability is $21/55$.  So the answer is $21+55=\boxed{076}$.

 




	\item (\textit{AIME II 2014}) A rectangle has sides of length $a$ and 36. A hinge is installed at each vertex of the rectangle, and at the midpoint of each side of length 36. The sides of length $a$ can be pressed toward each other keeping those two sides parallel so the rectangle becomes a convex hexagon as shown. When the figure is a hexagon with the sides of length $a$ parallel and separated by a distance of 24, the hexagon has the same area as the original rectangle. Find $a^2$. 


 


                    \begin{figure}[H]
                        \centering
                        \includegraphics[width=0.5\textwidth]{Problem/AIME_2014_II_3_Problem_Figure_1.png}
                    \end{figure}
                    


 
 \textbf{Answer: }720



\textbf{Solution 1}

When we squish the rectangle, the hexagon is composed of a rectangle and two isosceles triangles with side lengths 18, 18, and 24 as shown below.

 
                    \begin{figure}[H]
                        \centering
                        \includegraphics[width=0.5\textwidth]{Solution/AIME_2014_II_3_Solution_Figure_1.png}
                    \end{figure}
                    

 By Heron's Formula, the area of each isosceles triangle is $\sqrt{(30)(12)(12)(6)}=\sqrt{180\times 12^2}=72\sqrt{5}$. So the area of both is $144\sqrt{5}$. From the rectangle, our original area is $36a$. The area of the rectangle in the hexagon is $24a$. So we have 
 \[24a+144\sqrt{5}=36a\implies 12a=144\sqrt{5}\implies a=12\sqrt{5}\implies a^2=\boxed{720}.\]

 





\textbf{Solution 2}

Alternatively, use basic geometry. First, scale everything down by dividing everything by 6. Let $a/6=p$. Then, the dimensions of the central rectangle in the hexagon is p x 4, and the original rectangle is 6 x p. By Pythagorean theorem and splitting the end triangles of the hexagon into two right triangles, the altitude of the end triangles is $\sqrt{3^2-2^2}=\sqrt{5}$ given 2 as the base of the constituent right triangles. The two end triangles form a large rectangle of area $\sqrt{5}$ x $4$. Then, the area of the hexagon is $4p+4\sqrt{5}$, and the area of the rectangle is $6p$. Equating them, $p=2\sqrt{5}$. Multiply by scale factor of 6 and square it to get $36(20)= 720 \implies a^2=\boxed{720}$.

 




	\item (\textit{AIME II 2014}) The repeating decimals $0.abab\overline{ab}$ and $0.abcabc\overline{abc}$ satisfy 


 $0.abab\overline{ab}+0.abcabc\overline{abc}=\frac{33}{37},$


 where $a$, $b$, and $c$ are (not necessarily distinct) digits. Find the three digit number $abc$. 


 
 \textbf{Answer: }447



\textbf{Solution 1}

Notice repeating decimals can be written as the following:

 $0.\overline{ab}=\frac{10a+b}{99}$

 $0.\overline{abc}=\frac{100a+10b+c}{999}$

 where a,b,c are the digits. Now we plug this back into the original fraction:

 $\frac{10a+b}{99}+\frac{100a+10b+c}{999}=\frac{33}{37}$

 Multiply both sides by $999*99.$ This helps simplify the right side as well because $999=111*9=37*3*9$:

 $9990a+999b+9900a+990b+99c=33/37*37*3*9*99=33*3*9*99$

 Dividing both sides by $9$ and simplifying gives:

 $2210a+221b+11c=99^2=9801$

 At this point, seeing the $221$ factor common to both a and b is crucial to simplify. This is because taking $mod 221$ to both sides results in:

 $2210a+221b+11c  \equiv 9801 \mod 221 \iff 11c  \equiv 77 \mod 221$

 Notice that we arrived to the result $9801 \equiv 77 \mod 221$ by simply dividing $9801$ by $221$ and seeing $9801=44*221+77.$ Okay, now it's pretty clear to divide both sides by $11$ in the modular equation but we have to worry about $221$ being multiple of $11.$ Well, $220$ is a multiple of $11$ so clearly, $221$ couldn't be. Also, $221=13*17.$ Now finally we simplify and get:

 $c \equiv 7 \mod 221$

 But we know $c$ is between $0$ and $9$ because it is a digit, so $c$ must be $7.$ Now it is straightforward from here to find $a$ and $b$:

 $2210a+221b+11(7)=9801 \iff 221(10a+b)=9724 \iff 10a+b=44$

 and since a and b are both between $0$ and $9$, we have $a=b=4$. Finally we have the $3$ digit integer $\boxed{447}$

 





\textbf{Solution 2}

Note that $\frac{33}{37}=\frac{891}{999} = 0.\overline{891}$.  Also note that the period of $0.abab\overline{ab}+0.abcabc\overline{abc}$ is at most $6$.  Therefore, we only need to worry about the sum $0.ababab+ 0.abcabc$.  Adding the two, we get
 \[\begin{array}{ccccccc}&a&b&a&b&a&b\\ +&a&b&c&a&b&c\\ \hline &8&9&1&8&9&1\end{array}\]From this, we can see that $a=4$, $b=4$, and $c=7$, so our desired answer is $\boxed{447}$

 





\textbf{Solution 3}

Noting as above that $0.\overline{ab} = \frac{10a + b}{99}$ and $0.\overline{abc} = \frac{100a + 10b + c}{999}$, let $u = 10a + b$.Then
 \[\frac{u}{99} + \frac{10u + c}{999} = \frac{33}{37}\]

 
 \[\frac{u}{11} + \frac{10u + c}{111} = \frac{9\cdot 33}{37}\]

 
 \[\frac{221u + 11c}{11\cdot 111} = \frac{9\cdot 33}{37}\]

 
 \[221u + 11c = \frac{9\cdot 33\cdot 11\cdot 111}{37}\]

 
 \[221u + 11c = 9\cdot 33^2.\]

 Solving for $c$ gives

 
 \[c = 3\cdot 9\cdot 33 - \frac{221u}{11}\]

 
 \[c = 891 - \frac{221u}{11}\]

 Because $c$ must be integer, it follows that $u$ must be a multiple of $11$ (because $221$ clearly is not). Inspecting the equation, one finds that only $u = 44$ yields a digit $c, 7$. Thus $abc = 10u + c = \boxed{447}.$

 





\textbf{Solution 4}

We note as above that $0.\overline{ab} = \frac{10a + b}{99}$ and $0.\overline{abc} = \frac{100a + 10b + c}{999},$ so

 
 \[\frac{10a + b}{99} + \frac{100a + 10b + c}{999} = \frac{33}{37} = \frac{891}{999}.\]

 As $\frac{10a + b}{99}$ has a factor of $11$ in the denominator while the other two fractions don't, we need that $11$ to cancel, so $11$ divides $10a + b.$ It follows that $a = b,$ so $\frac{10a + b}{99} = \frac{11a}{99} = \frac{111a}{999},$ so

 
 \[\frac{111a}{999} + \frac{110a+c}{999} = \frac{891}{999}.\]

 Then $111a + 110a + c = 891,$ or $221a + c = 891.$ Thus $a = b = 4$ and $c = 7,$ so the three-digit integer $abc$ is $\boxed{447}.$

 




	\item (\textit{AIME II 2014}) Real numbers $r$ and $s$ are roots of $p(x)=x^3+ax+b$, and $r+4$ and $s-3$ are roots of $q(x)=x^3+ax+b+240$. Find the sum of all possible values of $|b|$.


 
 \textbf{Answer: }420



\textbf{Solution 1}

Because the coefficient of $x^2$ in both $p(x)$ and $q(x)$ is 0, the remaining root of $p(x)$ is $-(r+s)$, and the remaining root of $q(x)$ is $-(r+s+1)$. The coefficients of $x$ in $p(x)$ and $q(x)$ are both equal to $a$, and equating the two coefficients gives
 \[rs-(r+s)^2 = (r+4)(s-3)-(r+s+1)^2\]from which $s = \frac 12 (5r+13)$. Substitution should give $r = -5$ and $r = 1$, corresponding to $s = -6$ and $s = 9$, and $|b| = 330, 90$, for an answer of $\boxed{420}$.

 





\textbf{Solution 2}

Let $r$, $s$, and $-r-s$ be the roots of $p(x)$ (per Vieta's). Then $r^3 + ar + b = 0$ and similarly for $s$. Also,
 \[q(r+4) = (r+4)^3 + a(r+4) + b  + 240 = 12r^2 + 48r + 304 + 4a = 0\]

 Set up a similar equation for $s$:

 
 \[q(s-3) = (s-3)^3 + a(s-3) + b + 240 = -9s^2 + 27s + 213 - 3a = 0.\]

 Simplifying and adding the equations gives
 \begin{align}\tag{*} r^2 - s^2 + 4r + 3s + 49 &= 0 \end{align}Now, let's deal with the $ax$ terms. Plugging the roots $r$, $s$, and $-r-s$ into $p(x)$ yields a long polynomial, and plugging the roots $r+4$, $s-3$, and $-1-r-s$ into $q(x)$ yields another long polynomial. Equating the coefficients of $x$ in both polynomials, we get:
 \[rs + (-r-s)(r+s) = (r+4)(s-3) + (-r-s-1)(r+s+1),\]which eventually simplifies to
 \[s = \frac{13 + 5r}{2}.\]Substitution into (*) should give $r = -5$ and $r = 1$, corresponding to $s = -6$ and $s = 9$, and $|b| = 330, 90$, for an answer of $\boxed{420}$.

 





\textbf{Solution 3}

The roots of $p(x)$ are $r$, $s$, and $-r-s$ since they sum to $0$ by Vieta's Formula (coefficient of $x^2$ term is $0$). 

 Similarly, the roots of $q(x)$ are $r + 4$, $s - 3$, and $-r-s-1$, as they too sum to $0$. 

 Then:

 $a = rs+r(-r-s)+s(-r-s) = rs-(r+s)^2$ and $-b = rs(-r-s)$ from $p(x)$ and 

 $a=(r+4)(s-3)+(r+4)(-r-s-1)+(s-3)(-r-s-1) = (r+4)(s-3)-(r+s+1)^2$ and $-(b+240)=(r+4)(s-3)(-r-s-1)$ from $q(x)$. 

 From these equations, we can write that 
 \[rs-(r+s)^2 = (r+4)(s-3)-(r+s+1)^2 = a\]and simplifying gives 
 \[2s-5r-13=0 \Rightarrow s = \frac{5r+13}{2}.\]

 


 We now move to the other two equations regarding the product of the roots. We see that we can cancel a negative from both sides to get 
 \[rs(r+s) = b\] 
 \[(r+4)(s-3)(r+s+1)=b + 240.\] Subtracting the first equation from the second equation gives us $(r+4)(s-3)(r+s+1) - rs(r+s) = 240$. 

 Expanding, simplifying, substituting $s = \frac{5r+13}{2}$, and simplifying some more yields the simple quadratic $r^2 + 4r - 5 = 0$, so $r = -5, 1$. Then $s = -6, 9$. 

 Finally, we substitute back into $b=rs(r+s)$ to get $b = (-5)(-6)(-5-6) = -330$, or $b = (1)(9)(1 + 9) = 90$. 

 The answer is $|-330|+|90| = \boxed{420}$.

 





\textbf{Solution 4}

By Vieta's, we know that the sum of roots of $p(x)$ is $0$. Therefore,the roots of $p$ are $r, s, -r-s$. By similar reasoning, the roots of $q(x)$are $r + 4, s - 3, -r - s - 1$. Thus, $p(x) = (x - r)(x - s)(x + r + s)$and $q(x) = (x - r - 4)(x - s + 3)(x + r + s + 1)$.

 Since $p(x)$ and $q(x)$ have the same coefficient for $x$, we can go aheadand match those up to get
 \begin{align*}     rs - r(r + s) - s(r + s) &= (r + 4)(s - 3) - (r + 4)(r + s + 1) - (s - 3)(r + s + 1) \\     0 &= -13 - 5r + 2s \\     s &= \frac{5r + 13}{2} \end{align*}


 At this point, we can go ahead and compare the constant term in $p(x)$ and$q(x)$. Doing so is certainly valid, but we can actually do this another way. Notice that $p(s) = 0$. Therefore, $q(s) = 240$. If we plug that intoour expression, we get that
 \begin{align*}     q(s)  &= 3(s - r - 4)(r + 2s + 1) \\     240 &= 3(s - r - 4)(r + 2s + 1) \\     240 &= 3\left( \frac{3r + 5}{2} \right)(6r + 14) \\     80 &= (3r + 5)(3r + 7) \\     0 &= r^2 + 4r - 5 \end{align*}
This tells us that $(r, s) = (1, 9)$ or $(-5, -6)$. Since $-b$ is the product of the roots, we have that the two possibilities are $1 \cdot 9 \cdot (-10) = -90$and $(-5)(-6)11 = 330$. Adding the absolute values of these gives us$\boxed{420}$.

 




	\item (\textit{AIME II 2014}) Charles has two six-sided dice. One of the die is fair, and the other die is biased so that it comes up six with probability $\frac{2}{3}$ and each of the other five sides has probability $\frac{1}{15}$. Charles chooses one of the two dice at random and rolls it three times. Given that the first two rolls are both sixes, the probability that the third roll will also be a six is $\frac{p}{q}$, where $p$ and $q$ are relatively prime positive integers. Find $p+q$. 


 
 \textbf{Answer: }167



\textbf{Solution 1}

The probability that he rolls a six twice when using the fair die is $\frac{1}{6}\times \frac{1}{6}=\frac{1}{36}$.  The probability that he rolls a six twice using the biased die is $\frac{2}{3}\times \frac{2}{3}=\frac{4}{9}=\frac{16}{36}$.  Given that Charles rolled two sixes, we can see that it is $16$ times more likely that he chose the second die.  Therefore the probability that he is using the fair die is $\frac{1}{17}$, and the probability that he is using the biased die is $\frac{16}{17}$.  The probability of rolling a third six is

 
 \[\frac{1}{17}\times \frac{1}{6} + \frac{16}{17} \times \frac{2}{3} = \frac{1}{102}+\frac{32}{51}=\frac{65}{102}\]Therefore, our desired $p+q$ is $65+102= \boxed{167}$

 





\textbf{Solution 2 (Official Solution)}

This is an incredibly simple problem if one is familiar with conditional probability (SO BE FAMILIAR WITH IT)! The conditional probability that the third roll will be a six given that the first two rolls are sixes, is the conditional probability that Charles rolls three sixes given that his first two rolls are sixes. This is thus $\frac{\frac{1}{2}(\frac{2}{3})^3+\frac{1}{2}(\frac{1}{6})^3}{\frac{1}{2}(\frac{2}{3})^2+\frac{1}{2}(\frac{1}{6})^2}= \frac{65}{102}$. The answer is therefore $65+102= \boxed {167}$.

 Note: I have just found out that this is also the official solution. So I did not STEAL it, but as a coincidence, have come across the EXACT SAME SOLUTION. LIKE EXACTLY THE SAME. I AM SLIGHTLY FRIGHTENED. :/

 




	\item (\textit{AIME II 2014}) Let $f(x)=(x^2+3x+2)^{\cos(\pi x)}$. Find the sum of all positive integers $n$ for which 
$\left |\sum_{k=1}^n\log_{10}f(k)\right|=1.$


 
 \textbf{Answer: }021



\textbf{Solution 1}

First, let's split it into two cases to get rid of the absolute value sign

 $\left |\sum_{k=1}^n\log_{10}f(k)\right|=1 \iff \sum_{k=1}^n\log_{10}f(k)=1,-1$

 Now we simplify using product-sum logarithmic identites:

 $\log_{10}{f(1)}+\log_{10}{f(2)}+...+\log_{10}{f(n)}=\log_{10}{f(1)\cdot f(2) \cdot ... \cdot f(n)}=1,-1$

 


 Note that the exponent $\cos{\pi(x)}$ is either $-1$ if $x$ is odd or $1$ if $x$ is even. 

 Writing out the first terms we have 

 $\frac{1}{(2)(3)}(3)(4)\frac{1}{(4)(5)} \ldots$

 This product clearly telescopes (i.e. most terms cancel) and equals either $10$ or $\frac{1}{10}$. But the resulting term after telescoping depends on parity (odd/evenness), so we split it two cases, one where $n$ is odd and another where $n$ is even. 

 $\textbf{Case 1: Odd n}$

 For odd $n$, it telescopes to $\frac{1}{2(n+2)}$ where $n$ is clearly $3$. 

 $\textbf{Case 2: Even n}$

 For even $n$, it telescopes to $\frac{n+2}{2}$ where $18$ is the only possible $n$ value. Thus the answer is $\boxed{021}$

 





\textbf{Solution 2}

Note that $\cos(\pi x)$ is $-1$ when $x$ is odd and $1$ when $x$ is even.  Also note that $x^2+3x+2=(x+1)(x+2)$ for all $x$.  Therefore
 \[\log_{10}f(x)=\log_{10}(x+1)+\log_{10}(x+2)\ \ \ \text{if }x \text{ is even}\]
 \[\log_{10}f(x)=-\log_{10}(x+1)-\log_{10}(x+2)\ \ \ \text{if }x \text{ is odd}\]Because of this, $\sum_{k=1}^n\log_{10}f(k)$ is a telescoping series of logs, and we have
 \[\sum_{k=1}^n\log_{10}f(k)= \log_{10}(n+2)-\log_{10}2=\log_{10}\frac{n+2}{2}\ \ \ \text{if }n \text{ is even}\]
 \[\sum_{k=1}^n\log_{10}f(k)= -\log_{10}(n+2)-\log_{10}2=-\log_{10}2(n+2)\ \ \ \text{if }n \text{ is odd}\]Setting each of the above quantities to $1$ and $-1$ and solving for $n$, we get possible values of $n=3$ and $n=18$ so our desired answer is $3+18=\boxed{021}$

 




	\item (\textit{AIME II 2014}) Circle $C$ with radius 2 has diameter $\overline{AB}$. Circle $D$ is internally tangent to circle $C$ at $A$. Circle $E$ is internally tangent to circle $C$, externally tangent to circle $D$, and tangent to $\overline{AB}$. The radius of circle $D$ is three times the radius of circle $E$, and can be written in the form $\sqrt{m}-n$, where $m$ and $n$ are positive integers. Find $m+n$.


 
 \textbf{Answer: }254



\textbf{Solution 1}

<center>
                    \begin{figure}[H]
                        \centering
                        \includegraphics[width=0.5\textwidth]{Solution/AIME_2014_II_8_Solution_Figure_1.png}
                    \end{figure}
                    

 </center>Using the diagram above, let the radius of $D$ be $3r$, and the radius of $E$ be $r$. Then, $EF=r$, and $CE=2-r$, so the Pythagorean theorem in $\triangle CEF$ gives $CF=\sqrt{4-4r}$. Also, $CD=CA-AD=2-3r$, so 
 \[DF=DC+CF=2-3r+\sqrt{4-4r}.\] Noting that $DE=4r$, we can now use the Pythagorean theorem in $\triangle DEF$ to get 
 \[(2-3r+\sqrt{4-4r})^2+r^2=16r^2.\]

 Solving this quadratic is somewhat tedious, but the constant terms cancel, so the computation isn't terrible. Solving gives $3r=\sqrt{240}-14$ for a final answer of $\boxed{254}$.

 
                \begin{itemize}
                    \item Notice that C, E and the point of tangency to circle C for circle E will be collinear because C and E intersect the tangent line at a right angle, implying they must be on the same line.
                \end{itemize}


                





\textbf{Solution 2}

Consider a reflection of circle $E$ over diameter $\overline{AB}$. By symmetry, we now have three circles that are pairwise externally tangent and all internally tangent to a large circle. The small circles have radii $r$, $r$, and $3r$, and the big circle has radius $2$. 

 Descartes' Circle Theorem gives $\left(\frac{1}{r}+\frac{1}{r}+\frac{1}{3r}-\frac12\right)^2 = 2\left(\left(\frac{1}{r}\right)^2+\left(\frac{1}{r}\right)^2+\left(\frac{1}{3r}\right)^2+\left(-\frac12\right)^2\right)$

 Note that the big circle has curvature $-\frac12$ because it is internally tangent. Solving gives $3r=\sqrt{240}-14$ for a final answer of $\boxed{254}$.

 





\textbf{Solution 3}

We use the notation of Solution 1 for triangle $\triangle DEC$
 \[\sin \angle EDC = \frac {EF}{DE} = \frac {1}{4} \implies \cos \angle EDC =  \frac {\sqrt{15}}{4}.\]We use Cosine Law for $\triangle DEC$ and get:
 \[(4r)^2 +(2 – 3r)^2 – 2 \cdot 4r \cdot (2 – 3r) \cdot  \frac {\sqrt{15}}{4} = (2 – r)^2\].
 \[(24 + 6 \sqrt{15} ) r^2 = (8 + 4 \sqrt {15})r \implies 3r = 4 \sqrt{15} – 14 \implies \boxed{\textbf{254}}.\]\textbf{vladimir.shelomovskii@gmail.com, vvsss}

 





\textbf{Solution 4}

This problem can be very easily solved using Descartes' Circle Theorem. It states that if we have 4 circles that are all tangent with each other, $(k_1 + k_2 + k_3 + k_4)^{2} = 2(k_1^{2} + k_2^{2} + k_3^{2} + k_4^{2})$, where $k_i$ is the curvature of circle $i$, meaning $k_i = \dfrac{1}{r}$. When three of the circles are internally tangent to the fourth one, the fourth circle has a negative curvature. Suppose we reflect Circle $E$ over $\overline{AB}$. Now, we have our four circles to apply that theorem. First, lets scale our image down such that Circle $C$ has radius $1$, for ease of computation. Let the radius of Circle $D$ be $r$, so Circle $E$ has radius $\dfrac{r}{3}$. Then, we have that $(-1 + \dfrac{1}{r} + \dfrac{3}{r} + \dfrac{3}{r})^{2} = 2(1 + \dfrac{1}{r^{2}} + \dfrac{9}{r^{2}} + \dfrac{9}{r^{2}})$. This simplifies to $\dfrac{49}{r^{2}} - \dfrac{14}{r} + 1 = \dfrac{2r^{2} + 38}{r^{2}}$. Multiplying both sides by $r^{2}$, we get that $49 - 14r + r^{2} = 2r^{2} + 38$, or $r^2 + 14r - 11 = 0$. We get $r = -7 \pm 2\sqrt{15}$, but we want the positive solution, which is $r = 2\sqrt{15} - 7$. We have to rescale back up, so we get $r = 4\sqrt{15} - 14 = \sqrt{240} - 14$, so we get that our answer is $240 + 14 = \boxed{254}$.~Puck\_0

 





\textbf{Solution 5 (Heron's Formula)}

This solution focuses on the area of $\triangle DEC$. Because $EF$ is perpendicular to $AB$, it is an altitude of $\triangle DEC$. Therefore, we can express the area of $\triangle DEC$ as $\frac{1}{2}\cdot EF \cdot DC$. We can also express the area of $\triangle DEC$ using Heron's Formula. Let $r$ equal the radius of circle $E$. Then $DC$ = $CA - DA = 2-3r$. We also know that $CE = 2-r$ and $DE = 3r+r=4r$.The semi-perimeter of $\triangle DEC$ is $(DC+CE+DE)/2 = 2$. 

 Applying Heron's Formula, we get
 \[\sqrt{2(2-4r)(r)(3r)} = \sqrt{6r^2(2-4r)}.\]We set this equal to $\frac{1}{2}r(2-3r)$. 
 \[\sqrt{6r^2(2-4r)} = \frac{1}{2}r(2-3r).\]This simplifies to the quadratic equation $0 = 9r^2+84r-44$. Remember that we are solving for $3r$, which we will set equal to $x$. Then we now have the equation $0 = x^2 +28x - 44$. Applying the quadratic formula, we get $x=-14 \pm 4\sqrt{15}$. We want the positive solution, so we take $4\sqrt{15}-14 = \sqrt{240}-14$. Our answer is therefore $240 + 14 = \boxed{254}$.

 




	\item (\textit{AIME II 2014}) Ten chairs are arranged in a circle. Find the number of subsets of this set of chairs that contain at least three adjacent chairs.


 
 \textbf{Answer: }581



\textbf{Solution 1 (Casework)}

We know that a subset with less than $3$ chairs cannot contain $3$ adjacent chairs.  There are only $10$ sets of $3$ chairs so that they are all $3$ adjacent.  There are $10$ subsets of $4$ chairs where all $4$ are adjacent, and $10 \cdot 5$ or $50$ where there are only $3.$  If there are $5$ chairs, $10$ have all $5$ adjacent, $10 \cdot 4$ or $40$ have $4$ adjacent, and $10 \cdot {5\choose 2}$ or $100$ have $3$ adjacent.  With $6$ chairs in the subset, $10$ have all $6$ adjacent, $10(3)$ or $30$ have $5$ adjacent, $10 \cdot {4\choose2}$ or $60$ have $4$ adjacent, $\frac{10 \cdot 3}{2}$ or $15$ have $2$ groups of $3$ adjacent chairs, and $10 \cdot \left({5\choose2} - 3\right)$ or $70$ have $1$ group of $3$ adjacent chairs.  All possible subsets with more than $6$ chairs have at least $1$ group of $3$ adjacent chairs, so we add ${10\choose7}$ or $120$, ${10\choose8}$ or $45$, ${10\choose9}$ or $10$, and ${10\choose10}$ or $1.$  Adding, we get $10 + 10 + 50 + 10 + 40 + 100 + 10 + 30 + 60 + 15 + 70 + 120 + 45 + 10 + 1 = \boxed{581}.$

 





\textbf{Solution 2 (PIE)}

Starting with small cases, we see that four chairs give $4 + 1 = 5$, five chairs give $5 + 5 + 1 = 11$, and six chairs give $6 + 6 + 6 + 6 + 1 = 25.$ Thus, n chairs should give $n 2^{n-4} + 1$, as confirmed above. This claim can be verified by the principle of inclusion-exclusion: there are $n 2^{n-3}$ ways to arrange $3$ adjacent chairs, but then we subtract $n 2^{n-4}$ ways to arrange $4.$ Finally, we add $1$ to account for the full subset of chairs. Thus, for $n = 10$ we get a first count of $641.$

 However, we overcount cases in which there are two distinct groups of three or more chairs. We have $5$ cases for two groups of $3$ directly opposite each other, $5$ for two groups of four, $20$ for two groups of $3$ not symmetrically opposite, $20$ for a group of $3$ and a group of $4$, and $10$ for a group of $3$ and a group of $5.$ Thus, we have $641 - 60 = \boxed{581}$.

 





\textbf{Solution 3 (Complementary Counting and Recursion)}

It is possible to use recursion to count the complement. Number the chairs $1, 2, 3, ..., 10.$ If chair $1$ is not occupied, then we have a line of $9$ chairs such that there is no consecutive group of three. If chair $1$ is occupied, then we split into more cases. If chairs $2$ and $10$ are empty, then we have a line of $7.$ If chair $2$ is empty but chair $10$ is occupied, then we have a line of $6$ chairs (because chair $9$ cannot be occupied); this is similar to when chair $2$ is occupied and chair $10$ is empty. Finally, chairs $2$ and $10$ cannot be simultaneously occupied. Thus, we have reduced the problem down to computing $T_9 + T_7 + 2T_6$, where $T_n$ counts the ways to select a subset of chairs $\textit{in a line}$ from a group of n chairs such that there is no group of $3$ chairs in a row.

 Now, we notice that $T_n = T_{n-1} + T_{n-2} + T_{n-3}$ (representing the cases when the first, second, and/or third chair is unoccupied). Also, $T_0 = 1, T_1 = 2, T_2 = 4, T_3 = 7$, and hence $T_4 = 13, T_5 = 24, T_6 = 44, T_7 = 81, T_8 = 149, T_9 = 274$. Now we know the complement is $274 + 81 + 88 = 443$, and subtracting from $2^{10} = 1024$ gives $1024 - 443 = \boxed{581}$.

 





\textbf{Solution 4}

Let's calculate the complement. As mentioned in solution $3$, the number of ways to have a subset $n$ chairs in a line with no 3 consectuive chairs satisfies $T_n = T_{n-1} + T_{n-2} + T_{n-3}$. Setting $T_{1} = 2, T_{2} = 4$, and $T_{3} = 7$, we get that $T_{10} = 504$. Since this is in a line and not a circle, we must subtract the cases that would include 3 consecutive chairs if the endpoints of the line were put together. If chairs $1$, $2$ and $10$ are in the subset, that would not work. The same goes for if chairs $1$, $9$ and $10$ were in the subset. If chairs $1$, $2$ and $10$ are in the set then chair $3$ must not be in the set. However, chair $9$ could be in or not in the set because we only want to count what cases where no $3$ chairs are consecutive in the line but there would be consecutive chairs in a circle. If chair $9$ is not included, there are $T_{5} = 24$ ways. If chair $9$ is in the set then there are $T_{4} = 13$ ways. So we must subtract $2 \cdot (24 + 13) = 74$ However we are counting the case where chairs $1$, $2$, $9$ and $10$ are included twice. So we only have to subtract $74 - 13 = 61$. $504 - 61 = 443$. This entire time we were calculating the complement so $2^{10} - 443 = \fbox{581}$

 




	\item (\textit{AIME II 2014}) Let $z$ be a complex number with $|z|=2014$. Let $P$ be the polygon in the complex plane whose vertices are $z$ and every $w$ such that $\frac{1}{z+w}=\frac{1}{z}+\frac{1}{w}$. Then the area enclosed by $P$ can be written in the form $n\sqrt{3}$, where $n$ is an integer. Find the remainder when $n$ is divided by $1000$. 


 
 \textbf{Answer: }147



\textbf{Solution 1 (Roots of Unity)}

Notice that 
 \[\frac1{w+z} = \frac{w+z}{wz} \implies 0 = w^2 + wz + z^2 = \frac{w^3-z^3}{w-z}.\]Hence, $w=ze^{2\pi i/3},ze^{4\pi i/3}$, and $P$ is an equilateral triangle with circumradius $2014$. Then, 
 \[[P]=\frac{3}{2}\cdot 2014^2\cdot\sin\frac{\pi}3=3\cdot 1007^2\sqrt3,\]and the answer is $3\cdot 1007^2\equiv 3\cdot 7^2\equiv\boxed{147}\pmod{1000}$.

 





\textbf{Solution 2 (long but non-bashy)}

Note that the given equality reduces to

 
 \[\frac{1}{w+z} = \frac{w+z}{wz}\]
 \[wz = {(w+z)}^2\]
 \[w^2 + wz + z^2 = 0\]
 \[\frac{w^3 - z^3}{w-z} = 0\]
 \[w^3 = z^3, w \neq z\]

 Now, let $w = r_w e^{i \theta_w}$ and likewise for $z$. Consider circle $O$ with the origin as the center and radius 2014 on the complex plane. It is clear that $z$ must be one of the points on this circle, as $|z| = 2014$. 

 By DeMoivre's Theorem, the complex modulus of $w$ is cubed when $w$ is cubed. Thus $w$ must lie on $O$, since its the cube of its modulus, and thus its modulus, must be equal to $z$'s modulus.

 Again, by DeMoivre's Theorem, $\theta_w$ is tripled when $w$ is cubed and likewise for $z$. For $w$, $z$, and the origin to lie on the same line, $3 \theta_w$ must be some multiple of 360 degrees apart from $3 \theta_z$ , so $\theta_w$ must differ from $\theta_z$ by some multiple of 120 degrees.

 Now, without loss of generality, assume that $z$ is on the real axis. (The circle can be rotated to put $z$ in any other location.) Then there are precisely two possible distinct locations for $w$; one is obtained by going 120 degrees clockwise from $z$ about the circle and the other by moving the same amount counter-clockwise. Moving along the circle with any other multiple of 120 degrees in any direction will result in these three points.

 Let the two possible locations for $w$ be $W_1$ and $W_2$ and the location of $z$ be point $Z$. Note that by symmetry, $W_1W_2Z$ is equilateral, say, with side length $x$. We know that the circumradius of this equilateral triangle is $2014$, so using the formula $\frac{abc}{4R} = [ABC]$ and that the area of an equilateral triangle with side length $s$ is $\frac{s^2\sqrt{3}}{4}$, so we have

 
 \[\frac{x^3}{4R} = \frac{x^2\sqrt{3}}{4}\]
 \[x = R \sqrt{3}\]
 \[\frac{x^2\sqrt{3}}{4} = \frac{3R^2 \sqrt{3}}{4}\]

 Since we're concerned with the non-radical part of this expression and $R = 2014$,

 
 \[\frac{3R^2}{4} \equiv 3 \cdot 1007^2 \equiv 3 \cdot 7^2 \equiv \boxed{147} \pmod{1000}\]

 and we are done. $\blacksquare$

 





\textbf{Solution 3 (short)}

Assume $z = 2014$. Then
 \[\frac{1}{2014 + w} = \frac{1}{2014} + \frac{1}{w}\]

 
 \[2014w = w(2014 + w) + 2014(2014 + w)\]

 
 \[2014w = 2014w + w^2 + 2014^2 + 2014w\]

 
 \[0 = w^2 + 2014w + 2014^2\]

 
 \[w = \frac{-2014 \pm \sqrt{2014^2 - 4(2014^2)}}{2} = -1007 \pm 1007\sqrt{3}i\]

 Thus $P$ is an isosceles triangle with area $\frac{1}{2}(2014 - (-1007))(2\cdot 1007\sqrt{3}) = 3021\cdot 1007\sqrt{3}$ and $n \equiv 7\cdot 21\equiv \boxed{147} \pmod{1000}.$

 


 





\textbf{Solution 4}

I find that generally, the trick to these kinds of AIME problems is to interpret the problem geometrically, and that is just what I did here.  Looking at the initial equation, this seems like a difficult task, but rearranging yields a nicer equation:
 \[\frac1{z+w}=\frac1z+\frac1w\]
 \[\frac w{z+w}=\frac wz+1\]
 \[\frac w{z+w}=\frac{w+z}z\]
 \[\frac{w-0}{w-(-z)}=\frac{(-z)-w}{(-z)-0}\]We can interpret the difference of two complex numbers as a vector from one to the other, and we can interpret the quotient as a vector with an angle equal to the angle between the two vectors.  Therefore, after labeling the complex numbers with $W$ ($w$), $V$ ($-z$), and $O$ ($0$), we can interpret the above equation to mean that the $\angle OWV=\angle OVW$, and hence triangle $OWV$ is isosceles, so length $OW$ = $OV$.  Rearranging the equation
 \[\frac{w-0}{w-(-z)}=\frac{(-z)-w}{(-z)-0},\]we find that
 \[(w-0)((-z)-0)=-(w-(-z))^2.\]Taking the magnitude of both sides and using the fact that $OW=OV\implies |w-0|=|(-z)-0|$, we find that
 \[|w-0|^2=|w-(-z)|^2,\]so length $OW=VW$ and triangle $OWV$ is equilateral.  There are only two possible $W$'s for which $OWV$ is equilateral, lying on either side of $OV$.  After drawing these points on the circle of radius 2014 centered at the origin, it is easy to see that $z$ and the two $w$'s form an equilateral triangle (this can be verified by simple angle chasing).  Drawing a perpendicular bisector of one of the sides and using 30-60-90 triangles shows that the side length of this triangle is $2014\sqrt3$ and hence its area is $\frac{\sqrt3(2014\sqrt3)^2}4=\boxed{147}\sqrt3+1000k\sqrt3,$ for some integer $k$.

 


 




	\item (\textit{AIME II 2014}) In $\triangle RED$, $\measuredangle DRE=75^{\circ}$ and $\measuredangle RED=45^{\circ}$. $RD=1$. Let $M$ be the midpoint of segment $\overline{RD}$. Point $C$ lies on side $\overline{ED}$ such that $\overline{RC}\perp\overline{EM}$. Extend segment $\overline{DE}$ through $E$ to point $A$ such that $CA=AR$. Then $AE=\frac{a-\sqrt{b}}{c}$, where $a$ and $c$ are relatively prime positive integers, and $b$ is a positive integer. Find $a+b+c$.


 
 \textbf{Answer: }056



\textbf{Solution 1}

Let $P$ be the foot of the perpendicular from $A$ to $\overline{CR}$, so $\overline{AP}\parallel\overline{EM}$. Since triangle $ARC$ is isosceles, $P$ is the midpoint of $\overline{CR}$, and $\overline{PM}\parallel\overline{CD}$. Thus, $APME$ is a parallelogram and $AE = PM = \frac{CD}{2}$. We can then use coordinates. Let $O$ be the foot of altitude $RO$ and set $O$ as the origin. Now we notice special right triangles! In particular, $DO = \frac{1}{2}$ and $EO = RO = \frac{\sqrt{3}}{2}$, so $D\left(\frac{1}{2}, 0\right)$, $E\left(-\frac{\sqrt{3}}{2}, 0\right)$, and $R\left(0, \frac{\sqrt{3}}{2}\right).$ $M =$ midpoint$(D, R) = \left(\frac{1}{4}, \frac{\sqrt{3}}{4}\right)$ and the slope of $ME = \frac{\frac{\sqrt{3}}{4}}{\frac{1}{4} + \frac{\sqrt{3}}{2}} = \frac{\sqrt{3}}{1 + 2\sqrt{3}}$, so the slope of $RC = -\frac{1 + 2\sqrt{3}}{\sqrt{3}}.$ Instead of finding the equation of the line, we use the definition of slope: for every $CO = x$ to the left, we go $\frac{x(1 + 2\sqrt{3})}{\sqrt{3}} = \frac{\sqrt{3}}{2}$ up. Thus, $x = \frac{\frac{3}{2}}{1 + 2\sqrt{3}} = \frac{3}{4\sqrt{3} + 2} = \frac{3(4\sqrt{3} - 2)}{44} = \frac{6\sqrt{3} - 3}{22}.$ $DC = \frac{1}{2} - x = \frac{1}{2} - \frac{6\sqrt{3} - 3}{22} = \frac{14 - 6\sqrt{3}}{22}$, and $AE = \frac{7 - \sqrt{27}}{22}$, so the answer is $\boxed{056}$.

 
                    \begin{figure}[H]
                        \centering
                        \includegraphics[width=0.5\textwidth]{Solution/AIME_2014_II_11_Solution_Figure_1.png}
                    \end{figure}
                    

 





\textbf{Solution 2}

Let $MP = x.$ Meanwhile, since $\triangle R PM$ is similar to $\triangle RCD$ (angle, side, and side- $RP$ and $RC$ ratio), $CD$ must be 2$x$. Now, notice that $AE$ is $x$, because of the parallel segments $\overline A\overline E$ and $\overline P\overline M$.

 Now we just have to calculate $ED$. Using the Law of Sines, or perhaps using altitude $\overline R\overline O$, we get $ED = \frac{\sqrt{3}+1}{2}$. $CA=RA$, which equals $ED - x$

 Using Law of Sine in $\triangle RED$, we find $RE$ = $\frac{\sqrt{6}}{2}$. 

 We got the three sides of $\triangle AER$. Now using the Law of Cosines on $\angle AER$. There we can equate $x$ and solve for it. We got $AE=x=\frac{\sqrt{3}-1}{4\sqrt{3}+2}$. Then rationalize the denominator, we get $AE = \frac{7 - \sqrt{27}}{22}$.

 





\textbf{Solution 3}

Let $P$ be the foot of the perpendicular from $A$ to $\overline{CR}$, so $\overline{AP}\parallel\overline{EM}$. Since  $\triangle ARC$ is isosceles, $P$ is the midpoint of $\overline{CR}$, and by midpoint theorem $\overline{PM}\parallel\overline{CD}$. Thus, $APME$ is a parallelogram and therefore $AE = PM = \frac 12 CD$. 
                    \begin{figure}[H]
                        \centering
                        \includegraphics[width=0.5\textwidth]{Solution/AIME_2014_II_11_Solution_Figure_1.png}
                    \end{figure}
                    We can now use coordinates with $D(0,0)$ as origin and $DE$ along the $x$-axis. 

 Let $RD=4$ instead of $1$ (in the end we will scale down by $4$). Since $\angle D = 60^\circ$, we get $R(2,2\sqrt{3})$, and therefore $M(1, \sqrt{3})$. 

 We use sine-law in $\triangle RED$ to find the coordinates $E(2+2\sqrt{3}, 0)$:
 \[DE =4\cdot \frac{\sin 75^\circ}{\sin 45^\circ} = 4(\sin 30^\circ + \cos 30^\circ) = 2+2\sqrt{3}.\]Since slope$(ME)= -\sqrt{3}/(1+2\sqrt{3})$, and $RC\perp ME$, it follows that slope$(RC)=(1+2\sqrt{3})/\sqrt{3}$. If $C(c,0)$ then we have
 \[\frac{2\sqrt{3}}{2-c}=\frac{1+2\sqrt{3}}{\sqrt{3}}\qquad \Longrightarrow\qquad c=\frac{4\sqrt{3}-4}{1+2\sqrt{3}}\]Now $\frac 12 CD = \frac 12c =(2\sqrt{3}-2)/(1+2\sqrt{3})= \frac 1{11}(14-6\sqrt{3})$. 

 Scaling down by $4$, we get $AE=\frac 1{22}(7-3\sqrt{3})$, so our answer is $7+27+22=056$.

 





\textbf{Solution 4 (No intuition)}

Define the foot of the altitude from $R$ to $ED$ as $X.$ Suppose the median $EM$ is equal to $m,$ and the angle $\angle{MED}$ as $\theta.$ Using the Law of Cosines in triangle $EDM,$ we get $m^2=1+\frac{\sqrt{3}}{4}.$ By Law of Sines in $EMD,$ $\frac{m}{\sin{60^{\circ}}} = \frac{\frac{1}{2}}{\sin\theta}.$ This gives $\sin\theta = \frac{\sqrt{3}}{4m},$ and solving for $\cos\theta$ gives $\cos\theta = \sqrt{1-\frac{3}{16m^2}} = \frac{\sqrt{16m^2-3}}{4m}.$ Using the established value of $m^2,$ we have $\cos{\theta} = \frac{\sqrt{13+4\sqrt{3}}}{4m} = \frac{1+2\sqrt{3}}{4m}.$ Thus, we have $\tan\theta = \frac{\sqrt{3}}{1+2\sqrt{3}}.$ Noting that $\angle{XRC}=\theta$ as well, this gives $CX = RX \tan \theta = \frac{\sqrt{3}}{2} \cdot \frac{\sqrt{3}}{1+2\sqrt{3}} = \frac{3}{2+4\sqrt{3}}.$ Suppose $AE=n.$ Now, using Pythagorean in triangle $AXR$ gives 
 \[\left(n + \frac{\sqrt{3}}{2}\right)^2 + \frac{3}{4} = \left(n + \frac{\sqrt{3}}{2} + \frac{3}{2+4\sqrt{3}}\right)^2.\]Solving, we have
 \[\frac{3}{4} = \left(\frac{3}{2+4\sqrt{3}}\right) \left(2n + \sqrt{3} + \frac{3}{2+4\sqrt{3}}\right),\]
 \[2n + \sqrt{3} + \frac{3}{2+4\sqrt{3}} = \frac{1+2\sqrt{3}}{2},\]
 \[2n + \frac{3}{2+4\sqrt{3}} = \frac{1}{2}.\]Isolating $2n$ and rationalizing the denominator returns the final answer of $n=\frac{7-\sqrt{27}}{22},$ thus giving an answer of $\boxed{056}.$

 


 




	\item (\textit{AIME II 2014}) Suppose that the angles of $\triangle ABC$ satisfy $\cos(3A)+\cos(3B)+\cos(3C)=1$. Two sides of the triangle have lengths 10 and 13. There is a positive integer $m$ so that the maximum possible length for the remaining side of $\triangle ABC$ is $\sqrt{m}$. Find $m$. 


 
 \textbf{Answer: }399



\textbf{Solution 1}

Note that $\cos{3C}=-\cos{(3A+3B)}$. Thus, our expression is of the form $\cos{3A}+\cos{3B}-\cos{(3A+3B)}=1$. Let $\cos{3A}=x$ and $\cos{3B}=y$.

 Using the fact that $\cos(3A+3B)=\cos 3A\cos 3B-\sin 3A\sin 3B=xy-\sqrt{1-x^2}\sqrt{1-y^2}$, we get $x+y-xy+\sqrt{1-x^2}\sqrt{1-y^2}=1$, or $\sqrt{1-x^2}\sqrt{1-y^2}=xy-x-y+1=(x-1)(y-1)$. 

 Squaring both sides, we get $(1-x^2)(1-y^2) = [(x-1)(y-1)]^2$. Cancelling factors, $(1+x)(1+y) = (1-x)(1-y)$.

 
                \begin{itemize}
                    \item Notice here that we cancelled out one factor of (x-1) and (y-1), which implies that (x-1) and (y-1) were not 0. If indeed they were 0 though, we would have
                \end{itemize}


                $cos(3A)-1=0, cos(3A)=1$ 

 For this we could say that A must be 120 degrees for this to work. This is one case. The B case follows in the same way, where B must be equal to 120 degrees. This doesn't change the overall solution though, as then the other angles are irrelevant (this is the largest angle, implying that this will have the longest side and so we would want to have the 120 degreee angle opposite of the unknown side).

 Expanding, $1+x+y+xy=1-x-y+xy\rightarrow x+y=-x-y$. 

 Simplification leads to $x+y=0$.

 Therefore, $\cos(3C)=1$. So $\angle C$ could be $0^\circ$ or $120^\circ$. We eliminate $0^\circ$ and use law of cosines to get our answer: 

 
 \[m=10^2+13^2-2\cdot 10\cdot 13\cos\angle C\]
 \[\rightarrow m=269-260\cos 120^\circ=269-260\left(\text{-}\frac{1}{2}\right)\]
 \[\rightarrow m=269+130=399\]

 $\framebox{399}$

 
Note: We can get $x=-y$ from $(1+x)(1+y) = (1-x)(1-y)$, and plugging in $\cos{3A}+\cos{3B}=0$ directly gives $\cos(3C)=1$. Angle C being 120 degrees is also sufficient for $(1+x)(1+y) = (1-x)(1-y)$, so we can skip the casework part. -HappyByron

 





\textbf{Solution 2}

WLOG, let C be the largest angle in the triangle.

 As above, we can see that $\cos3A+\cos3B-\cos(3A+3B)=1$

 Expanding, we get

 $\cos3A+\cos3B-\cos3A\cos3B+\sin3A\sin3B=1$

 $\cos3A\cos3B-\cos3A-\cos3B+1=\sin3A\sin3B$

 $(\cos3A-1)(\cos3B-1)=\sin3A\sin3B$

 \textbf{CASE 1:} If $\sin 3A = 0$ or $\sin 3B = 0$,

 This implies one or both of A or B are 60 or 120.

 If one of A or B is 120, we have a contradiction, since C must be the largest angle.

 Otherwise, if one of A or B is 60, WLOG, assume A = 60, we would have $\cos(3B) + \cos(3C) = 2$, and thus, cos(3B) and cos(3C) both equal 1, implying $B = C = 120$, a contradiction to the fact that the sum of the angles of a triangle must be 180 degrees.

 \textbf{CASE 2:} If $\sin 3A \neq 0$ and $\sin 3B \neq 0$

 $\frac{\cos3A-1}{\sin3A}\cdot\frac{\cos3B-1}{\sin3B}=1$

 $\tan{\frac{3A}{2}}\tan{\frac{3B}{2}}=1$

 Note that $\tan{x}=\frac{1}{\tan(90-x)}$, or $\tan{x}\tan(90-x)=1$

 Thus $\frac{3A}{2}+\frac{3B}{2}=90$, or $A+B=60$.

 Now we know that $C=120$, so we can just use the Law of Cosines to get $\boxed{399}$

 -Alexlikemath

 





\textbf{Solution 3}


 \[\cos3A+\cos3B=1-\cos(3C)=1+\cos(3A+3B)\]
 \[2\cos\frac{3}{2}(A+B)\cos\frac{3}{2}(A-B)=2\cos^2\frac{3}{2}(A+B)\]If $\cos\frac{3}{2}(A+B) = 0$, then $\frac{3}{2}(A+B)=90$, $A+B=60$, so $C=120$; otherwise,
 \[2\cos\frac{3}{2}(A-B)=2 \cos\frac{3}{2}(A+B)\]
 \[\sin\frac{3}{2}A\sin\frac{3}{2}B=0\]so either $\sin\frac{3}{2}A=0$ or $\sin\frac{3}{2}B=0$, i.e., either $A=120$ or $B=120$. In all cases, one of the angles must be 120, which opposes the longest side. Final result follows. $\boxed{399}$

 -Mathdummy

 





\textbf{Solution 4(Fakesolve)}

Let $BC$ be the unknown side length. By Law of Cosines we have that $BC = \sqrt{269-260\cos{A}}$. We notice that $\cos{A}$ should be negative to optimize $BC$ so $A$ is between $90$ and $180$ degrees. We also know that the value inside the square root is an integer $m$, so $269-260\cos{A}$ should be an integer. We can then assume that $A$ is $120$ degrees so $\cos{A} = \frac{-1}{2}$. We do this because $120$ degrees is a "common" value and it makes the value inside the square root an integer. Plugging this into $269-260\cos{A}$ for $m$ we get that it is $\boxed{399}$. 

 -srisainandan6

 


 




	\item (\textit{AIME II 2014}) Ten adults enter a room, remove their shoes, and toss their shoes into a pile. Later, a child randomly pairs each left shoe with a right shoe without regard to which shoes belong together. The probability that for every positive integer $k<5$, no collection of $k$ pairs made by the child contains the shoes from exactly $k$ of the adults is $\frac{m}{n}$, where $m$ and $n$ are relatively prime positive integers. Find $m+n$. 


 
 \textbf{Answer: }028



\textbf{Solution}

Label the left shoes be $L_1,\dots, L_{10}$ and the right shoes $R_1,\dots, R_{10}$. Notice that there are $10!$ possible pairings.

 
Let a pairing be "bad" if it violates the stated condition. We would like a better condition to determine if a given pairing is bad.

 
Note that, in order to have a bad pairing, there must exist a collection of $k<5$ pairs that includes both the left and the right shoes of $k$ adults; in other words, it is bad if it is possible to pick $k$ pairs and properly redistribute all of its shoes to exactly $k$ people.

 
Thus, if a left shoe is a part of a bad collection, its corresponding right shoe must also be in the bad collection (and vice versa). To search for bad collections, we can start at an arbitrary right shoe (say $R_1$), check the left shoe it is paired with (say $L_i$), and from the previous observation, we know that $R_i$ must also be in the bad collection. Then we may check the left shoe paired with $R_i$, find its counterpart, check its left pair, find its counterpart, etc. until we have found $L_1$. We can imagine each right shoe "sending" us to another right shoe (via its paired left shoe) until we reach the starting right shoe, at which point we know that we have found a bad collection if we have done this less than $5$ times. 

 
Effectively we have just traversed a \textit{cycle.} (Note: This is the cycle notation of permutations.) The only condition for a bad pairing is that there is a cycle with length less than $5$; thus, we need to count pairings where every cycle has length at least $5$. This is only possible if there is a single cycle of length $10$ or two cycles of length $5$. 

 
The first case yields $9!$ working pairings. The second case yields $\frac{{10\choose 5}}{2}\cdot{4!}^2=\frac{10!}{2 \cdot {5!}^2} \cdot {4!}^2$ pairings. Therefore, taking these cases out of a total of $10!$, the probability is $\frac{1}{10}+\frac{1}{50} = \frac{3}{25}$, for an answer of $\boxed{028}$.

 


 MathProblemSolvingSkills.com

 


 \textbf{Article/Blog post solution} \href{https://www.commonsensiblemath.com/the-power-of-graphing-2014-aime-ii-p13/}{https://www.commonsensiblemath.com/the-power-of-graphing-2014-aime-ii-p13/}

 




	\item (\textit{AIME II 2014}) In $\triangle ABC$, $AB=10$, $\measuredangle A=30^{\circ}$, and $\measuredangle C=45^{\circ}$. Let $H$, $D$, and $M$ be points on line $\overline{BC}$ such that $AH\perp BC$, $\measuredangle BAD=\measuredangle CAD$, and $BM=CM$. Point $N$ is the midpoint of segment $HM$, and point $P$ is on ray $AD$ such that $PN\perp BC$. Then $AP^2=\frac{m}{n}$, where $m$ and $n$ are relatively prime positive integers. Find $m+n$.


 
 \textbf{Answer: }077



\textbf{Solution 1}

First, $\triangle DNP \sim \triangle DHA$, so 
 \[\frac{AP}{HN}=\frac{AH}{HD} = \sec \angle ADH = \sec 60^\circ= 2 \implies AP = 2HN = HM\]Now, by law of sines, $BC = 5\sqrt{2}$, and $HB=10\sin 15^\circ$. Thus, 
 \[AP^2 = HM^2 = (MB+BH)^2 = \left(\frac{5\sqrt{2}}{2}+\frac{10(\sqrt{6}-\sqrt{2})}{4} \right)^2 = \left(\frac{5\sqrt{6}}{2} \right)^2 = \frac{75}{2} \implies \boxed{077}\]

 -sillybone

 





\textbf{Solution 2}

Let us just drop the perpendicular from $B$ to $AC$ and label the point of intersection $O$. We will use this point later in the problem.As we can see, $M$ is the midpoint of $BC$ and $N$ is the midpoint of $HM$,$AHC$ is a $45-45-90$ triangle, so $\angle{HAB}=15^\circ$.$AHD$ is $30-60-90$ triangle.

 $AH$ and $PN$ are parallel lines so $PND$ is a $30-60-90$ triangle also.Then if we use those informations we get $AD=2HD$ and $PD=2ND$ and $AP=AD-PD=2HD-2ND=2HN$ or $AP=2HN=HM$.Now we know  that $HM=AP$, we can find for $HM$ which is simpler to find.We can use point $B$ to split it up as $HM=HB+BM$,We can chase those lengths and we would get $AB=10$, so $OB=5$, so $BC=5\sqrt{2}$, so $BM=\dfrac{1}{2} \cdot BC=\dfrac{5\sqrt{2}}{2}$We can also use Law of Sines:
 \[\frac{BC}{AB}=\frac{\sin\angle A}{\sin\angle C}\]
 \[\frac{BC}{10}=\frac{\frac{1}{2}}{\frac{\sqrt{2}}{2}}\implies BC=5\sqrt{2}\]Then using right triangle $AHB$, we have $HB=10 \sin 15^\circ$So $HB=10 \sin 15^\circ=\dfrac{5(\sqrt{6}-\sqrt{2})}{2}$.And we know that $AP = HM = HB + BM = \frac{5(\sqrt6-\sqrt2)}{2} + \frac{5\sqrt2}{2} = \frac{5\sqrt6}{2}$.Finally, $(AP)^2=\dfrac{150}{4}=\dfrac{75}{2}$. So our final answer is $75+2=77$. $m+n=\boxed{077}$

 -Gamjawon

 -edited by srisainandan6 to clarify and correct a small mistake

 





\textbf{Solution 3}

Here's a solution that doesn't need $\sin 15^\circ$.

 As above, get to $AP=HM$.  As in the figure, let $O$ be the foot of the perpendicular from $B$ to $AC$.  Then $BCO$ is a 45-45-90 triangle, and $ABO$ is a 30-60-90 triangle.  So $BO=5$ and $AO=5\sqrt{3}$; also, $CO=5$, $BC=5\sqrt2$, and $MC=\dfrac{BC}{2}=5\dfrac{\sqrt2}{2}$.  But $MO$ and $AH$ are parallel, both being orthogonal to $BC$.  Therefore $MH:AO=MC:CO$, or $MH=\dfrac{5\sqrt3}{\sqrt2}$, and we're done.

 





\textbf{Solution 4}

Break our diagram into 2 special right triangle by dropping an altitude from $B$ to $AC$ we then get that 
 \[AC=5+5\sqrt{3}, BC=5\sqrt{2}.\]Since $\triangle{HCA}$ is a 45-45-90,

 
 \[HC=\frac{5\sqrt2+5\sqrt6}{2}\] $MC=\frac{BC}{2},$ 
 \[HM=\frac{5\sqrt6}{2}\] 
 \[HN=\frac{5\sqrt6}{4}\] We know that $\triangle{AHD}\simeq \triangle{PND}$ and are 30-60-90.  Thus, 
 \[AP=2 \cdot HN=\frac{5\sqrt6}{2}.\]

 $(AP)^2=\dfrac{150}{4}=\dfrac{75}{2}$. So our final answer is $75+2=\boxed{077}$.

 





\textbf{Solution 5}


                    \begin{figure}[H]
                        \centering
                        \includegraphics[width=0.5\textwidth]{Solution/AIME_2014_II_14_Solution_Figure_1.png}
                    \end{figure}
                    Draw the $45-45-90 \triangle AHC$. Now, take the perpendicular bisector of $BC$ to intersect the circumcircle of $\triangle ABC$ and $AC$ at $F, L, G$ as shown, and denote $O$ to be the circumcenter of $\triangle ABC$. It is not difficult to see by angle chasing that $AHBGO$ is cyclic, namely with diameter $AB$. Then, by symmetry, $EH = HB$ and as $HB, OG$ are both subtended by equal arcs they are equal. Hence, $EH = GO$. Now, draw line $HL$ and intersect it at $AC$ at point $K$ in the diagram. It is not hard to use angle chase to arrive at $AEOL$ a parallelogram, and from our length condition derived earlier, $AL \parallel HG$. From here, it is clear that $AK = KG$; that is, $P$ is just the intersection of the perpendicular from $K$ down to $BC$ and $AD$! After this point, note that $AP = PF$. It is easily derived that the circumradius of $\triangle ABC$ is $\frac{10}{\sqrt{2}}$. Now, $APO$ is a $30-60-90$ triangle, and from here it is easy to arrive at the final answer of $\boxed{077}$. ~awang11's sol

 





\textbf{Solution 6}


                    \begin{figure}[H]
                        \centering
                        \includegraphics[width=0.5\textwidth]{Solution/AIME_2014_II_14_Solution_Figure_1.png}
                    \end{figure}
                    Let $BO \perp AC, O \in AC.$

 Let $ME \perp BC, E \in AD.$

 $MB = MC, \angle C = 45^\circ \implies$ points $M, E, O$ are collinear.

 $HN = NM, AH||NP||ME \implies AP = PE.$

 In $\triangle ABO \hspace{10mm} \angle A = 30^\circ \implies AO = AB \cos 60^\circ  = 5 \sqrt{3}.$

 In $\triangle AEO \hspace{10mm} \angle A = 15^\circ, \angle O = 90^\circ + 45^\circ = 135^\circ  \implies$
 \[\angle AEO = 30^\circ \implies\]
 \[AE = AO \frac {\sin 135^\circ}{\sin 30^\circ}  = 5 \sqrt{3} \cdot \sqrt{2} = 5 \sqrt{6} \implies\]
 \[AP = 5 \sqrt {\frac {3}{2}} \implies AP^2 = \frac {75}{2} \implies  \boxed{\textbf{077}}.\]\textbf{vladimir.shelomovskii@gmail.com, vvsss}

 


 




	\item (\textit{AIME II 2014}) For any integer $k\geq 1$, let $p(k)$ be the smallest prime which does not divide $k$. Define the integer function $X(k)$ to be the product of all primes less than $p(k)$ if $p(k)>2$, and $X(k)=1$ if $p(k)=2$. Let $\{x_n\}$ be the sequence defined by $x_0=1$, and $x_{n+1}X(x_n)=x_np(x_n)$ for $n\geq 0$. Find the smallest positive integer $t$ such that $x_t=2090$. 


 
 \textbf{Answer: }149



\textbf{Solution}

Note that for any $x_i$, for any prime $p$, $p^2 \nmid x_i$.  This provides motivation to translate $x_i$ into a binary sequence $y_i$. 

 Let the prime factorization of $x_i$ be written as $p_{a_1} \cdot p_{a_2} \cdot p_{a_3} \cdots$, where $p_i$ is the $i$th prime number. Then, for every $p_{a_k}$ in the prime factorization of $x_i$, place a $1$ in the $a_k$th digit of $y_i$. This will result in the conversion $x_1 = 2, x_{2} = 3, x_{3} = 2 * 3 = 6, \cdots$.

 Multiplication for the sequence $x_i$ will translate to addition for the sequence $y_i$. Thus, we see that $x_{n+1}X(x_n) = x_np(x_n)$ translates into $y_{n+1} = y_n+1$. Since $x_0=1$, and $y_0=0$, $x_i$ corresponds to $y_i$, which is $i$ in binary. Since $x_{10010101_2} = 2 \cdot 5 \cdot 11 \cdot 19 = 2090$, $t = 10010101_2$ = $\boxed{149}$.

 





\textbf{Solution 2 DO NOT DO (Painful Bash)}

We go through the terms and look for a pattern. We find that

 $x_0 = 1$ $x_8 = 7$

 $x_1 = 2$ $x_9 = 14$

 $x_2 = 3$ $x_{10} = 21$

 $x_3 = 6$ $x_{11} = 42$

 $x_4 = 5$ $x_{12} = 35$

 $x_5 = 10$ $x_{13} = 70$

 $x_6 = 15$ $x_{14} = 105$

 $x_7 = 30$ $x_{15} = 210$

 Commit to the bash. Eventually, you will receive that $x_{149} = 2090$, so $\boxed{149}$ is the answer. Trust me, this is worth the 10 (20*) index points.

 $\textbf{-RootThreeOverTwo}$

 NOTE: this only works for this problem because the answer is a low number like 149, and the numbers are small. DO NOT try this in a real problem unless it is your last one, as this is very risky. ~eqb5000

 





\textbf{Solution 3 (induction)}

Let $P_k$ denote the $k$th prime.Lemma: $x_{n+2^{k-1}} = P_k \cdot x_{n}$ for all $0 \leq n \leq 2^{k-1} - 1.$

 $\mathbf{\mathrm{Proof:}}$

 We can prove this using induction. Assume that $x_{2^{k-1}-1} = \prod_{j=1}^{k-1} P_j.$ Then, using the given recursion $x_{k+1} = \frac{x_np(x_n)}{X(x_n)}$, we would “start fresh” for $x_{2^{k-1}} = P_k.$ It is then easy to see that then $\frac{x_n}{P_k}$ just cycles through the previous $x_{2^{k-1}}$ terms of $\{ x_n \},$ since the recursion process is the same and $P_k$ being a factor of $x_n$ is not affected until $n = 2 \cdot {2^{k-1}} = 2^k,$ when given our assumption $x_{2^k - 1} = \prod_{j=1}^{k} P_j$ and $n = 2^k$ is now the least $n$ such that 
 \[P_{k+1} = p(x_{2^{n-1}}),\] in which $P_a = p(x_n)$ where $a > k$ is the only way that the aforementioned cycle would be affected. Specifically, by cancellation according to our recursion, $x_{2^k} = P_{k+1},$ and the values of $x_n$ just starts cycling through the previous $x_{2^k}$ terms again until $x_{2^{k+1}}$ when a new prime shows up in the prime factorization of $x_n,$ when it starts cycling again, and so on. Using our base cases of $x_0$ and $x_1,$ our induction is complete. Now, it is easy to see that $2090 = 2 \cdot 5 \cdot 11 \cdot 19 = P_1 \cdot P_3 \cdot P_5 \cdot P_8,$ and by Lemma \#1, the least positive integer $n$ such that $19 | x_n$ is $2^7.$ By repeatedly applying our obtained recursion from Lemma \#1, it is easy to see that our answer is just $2^7 + 2^4 + 2^2 + 2^0,$ or $10010101_2 = \boxed{149}.$

 -fidgetboss\_4000

 


 


 




\end{enumerate}